\documentclass[../../main.tex]{subfiles}
\begin{document}


{\bf [解]}
与えられた関数の絶対値の中身を
\begin{align}
 f(x)=e^x-ax \label{eq:1} 
\end{align}
とおくと一階微分は
\begin{align}
 f'(x)=e^x-a \label{eq:2}
\end{align}
である．$[0,1]$では$1 \le e^{x} \le e$だから，$a$の値に応じて$f(x)$の挙動は以下のように変化する．


\subparagraph{$1^\circ\ a\le1$ の時}

この時は$e^{x} \ge a$より $f'(x)\ge0$ となって $f(x)$ は単調増加である．
$f(0)=1,f(1)=e-a$ だから，
\begin{align}
\max|f(x)|=2 \iff e-a=2 \iff a=e-2 \label{eq:3}
\end{align}
でこれは $a\le1$ をみたすから十分である．

\subparagraph{$2^\circ\ 1\le a\le e$ の時}

$f'(x)=0$ なる $x$ が $0\le x\le1$ にただ1つ存在するからこれを$\alpha$とおく．
$f(x)$の増減表は下表のようになる．

\begin{table}[htb]
\centering
\caption{$f(x)$ の増減表 ($1 \le a \le e$)}
\begin{tabular}{c|ccccc}
$x$ & $0$ & $\cdots$ & $\alpha$ & $\cdots$ & $1$ \\
\hline
$f'(x)$ & & $-$ & $0$ & $+$ & \\
\hline
$f(x)$ & $1$ & $\searrow$ & $e^\alpha(1-\alpha)$ & $\nearrow$ & $e-a$
\end{tabular}
\end{table}

ここで$f'(\alpha)=0$ から，$e^\alpha=a$ で $1\le a\le e$ から $0\le\alpha\le1$であることに注意する．
$\max|f(x)|=2$ となる時を考えると，$|f(x)|$の最大値は端点または極値でとる．
しかし$|f(1)|=|e-a|<2$より端点はいずれも条件を満たさないから増減表から
\begin{align}
 f(\alpha) = -2
\end{align}
となるしかないが，
\begin{align}
  f(\alpha)=e^\alpha-a\alpha=e^\alpha(1-\alpha)>0
\end{align}
から矛盾する．よってこの時は条件を満たさない．

\subparagraph{$3^\circ\ e\le a$ の時}

この時は$e^x \le a$ より $f'(x)\le0$ であって $f(x)$ は単調減少である．
$f(0)=1,f(1)=e-a \le 0$ だから
\begin{align}
\max|f(x)|=2\iff e-a=-2\iff a=e+2 \label{eq:4}
\end{align}
でこれは $a\ge e$ をみたすから十分である．
\casesend

以上\cref{eq:3,eq:4}から，もとめる$a$は
\begin{align}
 a=e\pm2
\end{align}
である．

\newpage

\end{document}
